\documentclass[10pt,a4paper]{amsart}
\usepackage[margin=19mm]{geometry}
\usepackage{amsmath,amssymb,mathtools}
\usepackage[scr=rsfs,cal=euler]{mathalfa}
\usepackage{xfrac}
\usepackage[unicode,hidelinks,bookmarks=false]{hyperref}
\newcommand{\ie}{i.e.\ }
\newcommand{\cf}{cf.\ }
\newcommand{\resp}{respectively\ }
\newcommand{\Subsection}{\subsection}
\providecommand{\nobreakdash}{\nobreak\hbox{-}}
\setlength{\emergencystretch}{2em}
\multlinegap0pt
\usepackage{amssymb}
\usepackage{xcolor}
\usepackage{tikz}
\usepackage{algorithm}\usepackage{algpseudocode}
\newtheorem{theorem}{Theorem}
\newcommand{\sm}{{\mathcal{SM}}}
\newcommand{\zs}{{\mathbb S}}
\title{Complexes of strong discrete Morse matchings}
\author{Kevin Knudson, Abigail Owens-White}
\date{Source: arXiv:2609.06844v1 (CC BY 4.0)}
\begin{document}
\maketitle

\noindent\emph{Source and licence.} Complexes of strong discrete Morse matchings. Version arXiv:2609.06844v1 (CC BY 4.0), distributed by arXiv under the Creative Commons Attribution 4.0 International licence, http://creativecommons.org/licenses/by/4.0/, as recorded in the arXiv licence metadata for this exact version and shown on the abstract page https://arxiv.org/abs/2609.06844v1. Full licence notice: \texttt{licenses/arxiv-2609.06844-CC-BY-4.0.txt}.\par\smallskip

\noindent\textit{Editorial scope.} Counted unit: the article's theorem startheorem (source0.tex lines 647--649). The article's complete original proof of it is delivered with it (source0.tex lines 651--661, one \texttt{proof} environment of 11 lines). No mathematics is written, repaired or re-derived: the statement and the proof are the article's own lines, in the article's own notation, and no retained line is substituted or re-flowed.  No further source-local context is needed: every cross-reference of every retained line resolves inside the two delivered spans.  Nothing was omitted from any retained span, so the package carries no omission record.  No retained line contains a citation, so the wrapper carries no bibliography and no bibliography entry.  No retained line contains a figure environment, an image-inclusion command or drawing code, so no image is delivered and the package declares no asset dependency.  Editorial formatting only, in addition: an AMS \texttt{amsart} preamble that restates the article's own declarations and macros used by the retained lines, namely the environments \texttt{theorem} on separate counters, together with the standard packages those lines need.  The retained environments print in this document as Theorem 1 in order of appearance, whereas the article prints the same items with the numbering of its own full document.  The complete native source of the article as distributed in the arXiv e-print for this exact version is bundled unchanged as \texttt{sources/209578/source0.tex}.
\begin{theorem}\label{startheorem}
    For all $n\ge 2$, $\sm(S_n)$ is contractible.
\end{theorem}

\begin{proof}
    Note that $S_1=P_1$ and hence $\sm(S_1)\simeq \zs^0$. Now assume $n\ge 2$. The complex $\sm(P_n)$ has $2n$ vertices corresponding to the $2n$ possible vertex-edge pairs in $S_n$. Order these lexicographically--$01,02,\dots ,0n,10,20,\dots ,n0$--and then number them $1,2,\dots ,2n$.  
    
    If we label vertex $v_0$ with $0$, then for any labeling of the other $n$ vertices, we have $\textrm{lk}^\downarrow(v_i) = v_0$. We therefore obtain the maximal matching in which every vertex $v_i$ is paired with the edge joining it to $v_0$ and $v_0$ is left critical. The corresponding facet of $\sm(S_n)$ is $$[n+1,n+2,\dots ,2n].$$ 

    Now suppose that $v_i$ is labeled with $0$ and $v_0$ is labeled with $1$. For any ordering of the other $n-1$ vertices the lower link will consist of $v_0$ and so we obtain the maximal matchings $$[i,n+1,n+2,\dots ,\widehat{n+i},\dots ,2n]$$ for $1\le i\le n$.
    
    Finally consider labelings where $v_0$ is assigned value $k\ge 2$. The $k$ vertices labeled with values $0,\dots ,k-1$ are unpaired, and when $v_0$ gets its label, these vertices form $\textrm{lk}^\downarrow(v_0)$ and hence $v_0$ is not paired with any edge. If $v_j$ is any remaining vertex, then $\textrm{lk}^\downarrow(v_j) = \{v_0\}$. If $v_{i_1},\dots ,v_{i_{n-k}}$ are these vertices, then we obtain the matching $$[n+i_1,n+i_2,\dots ,n+{i_{n-k}}],$$ which is a face of the facet $[n+1,n+2,\dots ,2n]$ above. It follows that $\sm(S_n)$ is pure with the $n+1$ $(n-1)$-simplices listed above as facets.

    The complex $\sm(S_n)$ is shellable, with shelling order $$[n+1,n+2,\dots ,2n],[1,n+2,\dots ,2n],[2,n+1,n+3,\dots,2n],\dots ,[n,n+1,\dots ,2n-1].$$ There are no generating simplices since for each $i$, $F_i\cap (\bigcup_{j=1}^{i-1}F_j)$ is not the entire boundary of $F_i$ (the initial vertex of each $F_i$ does not appear in any of its predecessors) and so $\sm(S_n)$ is contractible.
\end{proof}

\end{document}
